% Apresentacao Beamer -- Poco 861 Etapa 1 (campo Lapa)
% Compilar: make beamer1

\documentclass[aspectratio=169,10pt]{beamer}

\usepackage[T1]{fontenc}
\usepackage[utf8]{inputenc}
\usepackage[brazil]{babel}
% DRP template uses Calibri; Carlito is the TeX metric-compatible substitute.
\usepackage[sfdefault]{carlito}
\usefonttheme{professionalfonts}
\usepackage{amsmath,amssymb}
\usepackage{booktabs}
\usepackage{array}
\usepackage{graphicx}
\usepackage{tikz}
\usetikzlibrary{positioning, arrows.meta, calc}

\graphicspath{%
  {figures/}%
  {../data/processed/ml_runs/diagnostics_861/figures/}%
  {../data/processed/ml_runs/well_profile/phi_lab/rf/}%
  {../data/processed/ml_runs/well_profile/phi_lab/compare/figures/}%
  {../data/processed/ml_runs/well_profile/phi_lab/alternatives/}%
  {../data/processed/ml_runs/well_profile/hfu/classifier/}%
  {../data/processed/ml_runs/ct_plugs/by_target/FZI_lab/wireline_only/}%
}

\usetheme{Madrid}
\usecolortheme{whale}
\definecolor{po861blue}{RGB}{0,74,127}
\definecolor{po861green}{RGB}{0,128,96}
\setbeamercolor{structure}{fg=po861blue}
\setbeamercolor{frametitle}{bg=po861blue,fg=white}
\setbeamercolor{title}{fg=white,bg=po861blue}
\setbeamercolor{subtitle}{fg=po861blue}
\setbeamertemplate{navigation symbols}{}
\setbeamertemplate{blocks}[rounded][shadow=false]

\setbeamersize{text margin left=9mm, text margin right=9mm}
\setbeamerfont{normal text}{size=\small}
\setbeamerfont{itemize/enumerate body}{size=\small}
\setbeamerfont{itemize/enumerate subbody}{size=\footnotesize}
\setbeamerfont{block title}{size=\small}
\setbeamerfont{block body}{size=\footnotesize}
\setbeamerfont{frametitle}{size=\large}
\setbeamerfont{title}{size=\normalsize}
\setbeamerfont{subtitle}{size=\small}

\setlength{\parskip}{2pt}
\setlength{\itemsep}{2pt}

\setbeamertemplate{footline}{
  \hfill
  \usebeamercolor[fg]{structure}
  \insertframenumber/\inserttotalframenumber\kern1em\vspace{2pt}
}

\newcommand{\philab}{\ensuremath{\phi_{\mathrm{lab}}}}
\newcommand{\fzilab}{\ensuremath{\mathrm{FZI}_{\mathrm{lab}}}}
\newcommand{\klab}{\ensuremath{k_{\mathrm{lab}}}}
\newcommand{\phiND}{\ensuremath{\phi_{\mathrm{ND}}}}

\newcommand{\slidefig}[2][0.58]{%
  \includegraphics[width=\linewidth,height=#1\textheight,keepaspectratio]{#2}%
}
\newcommand{\slidecaption}[1]{\\[-0.05cm]{\scriptsize #1}}
\newcommand{\oofcv}{validação \emph{out-of-fold} (OOF) em 3 blocos de profundidade ($n=87$)}
\newcommand{\rSqOOF}{$R^2$ OOF médio (3 blocos)}

\title{Poço~861 (campo Lapa): integração de dados\\e modelos preditivos}
\subtitle{Etapa~1 --- limites do perfil antes da física de rochas}
\author{Romulo Brito da Silva}
\institute{UENF / LENEP}
\date{Junho de 2026}

\begin{document}

% ------------------------------------------------------------------
\begin{frame}[plain]
  \titlepage
  \vspace{0.25cm}
  \begin{center}
    \footnotesize
    Intervalo 5205,91--5233,72~m $\cdot$ campo Lapa (pré-sal)\\
    87 amostras perfil+laboratório $\cdot$ 10 plugs com microtomografia (microCT)\\[0.12cm]
    Etapa diagnóstica: ainda não há $V_p$/$V_s$ teóricos
  \end{center}
\end{frame}

% ------------------------------------------------------------------
\begin{frame}{Roteiro}
  \footnotesize
  \begin{enumerate}
    \item Intervalo do poço~861 e montagem do \emph{perfil enriquecido}.
    \item Protocolo de validação: 87 linhas (blocos OOF) e 10 plugs (LOPO).
    \item Correlações perfil--laboratório e teste \emph{oracle} de FZI.
    \item Resultados do Random Forest: \philab, FZI, $k$ e HFU; comparacao
          \phiND{} direto versus RF.
    \item Encaminhamento para DEM/SC (Etapa~2).
  \end{enumerate}
\end{frame}

% ------------------------------------------------------------------
\section{Contexto}

\begin{frame}{O intervalo do poço~861 antes de modelar $V_p$}
  \footnotesize
  Carbonatos pré-sal do campo Lapa, $\sim$28~m de coluna (5206--5234~m).
  \philab, \klab{} e \fzilab{} variam bastante ao longo do trecho
  (painéis: 87 medidas de laboratório).
  \vspace{0.08cm}
  \begin{columns}[T,onlytextwidth]
    \column{0.32\textwidth}
    \begin{itemize}
      \item \textbf{Três escalas:} perfil (87 linhas), laboratório,
            micro-CT (10 plugs).
      \item \textbf{Alvos testados:} \philab, \fzilab, \klab, HFU.
      \item Nesta etapa \textbf{não} estimamos $V_p$ nem $V_s$ diretamente.
    \end{itemize}
    \vspace{0.08cm}
    {\tiny
    \textbf{Médias por bloco de validação OOF}
    \vspace{0.02cm}
    \begin{tabular}{@{}ccrrr@{}}
      \toprule
      Bl. & Prof.\ (m) & $\overline{\phi}$ & $\overline{k}$ & $\overline{\mathrm{FZI}}$ \\
      & & (pu) & (mD) & \\
      \midrule
      0 & 5206--5214 & 0,092 & 3,6 & 1,39 \\
      1 & 5215--5223 & 0,115 & 9,5 & 1,25 \\
      2 & 5225--5234 & 0,145 & 54,6 & 1,87 \\
      \bottomrule
    \end{tabular}}
    \column{0.66\textwidth}
    \centering
    \slidefig[0.58]{mogno_lab_properties_vs_depth}
    \slidecaption{Esquerda: \philab; centro: \klab{} (mD, escala log);
    direita: \fzilab{} vs.\ profundidade ($n=87$ pontos de laboratório).}
  \end{columns}
  \vspace{0.04cm}
  \begin{block}{Objetivo da Etapa~1}
    \footnotesize Verificar se GR, densidade, resistividade e porosidades de
    perfil prevêem \philab, \fzilab, \klab{} e HFU medidos em laboratório,
    com validação por blocos de profundidade --- antes de partir para $V_p$
    e $V_s$.
  \end{block}
\end{frame}

% ------------------------------------------------------------------
\section{Dados}

\begin{frame}{Perfil enriquecido: como montamos a tabela}
  \scriptsize
  Alinhamos perfil, laboratório e microtomografia na mesma profundidade.
  O laboratório entra só como alvo ($y$), nunca como entrada ($X$).
  \vspace{0.05cm}
  \begin{columns}[T,onlytextwidth]
    \column{0.48\textwidth}
    \begin{itemize}
      \item \textbf{Coluna de perfil} ($n=87$): curvas wireline e, em cada
            ponto, \philab, \klab, \fzilab, HFU de laboratório.
      \item \textbf{Dez plugs microCT}: porosidade de imagem, textura e
            composição mineral em escala milimétrica.
      \item \textbf{Casamento} ($\pm$0,5~m): cada plug vai para a linha de
            perfil mais próxima; oito profundidades recebem descritores de CT.
      \item \textbf{Dois usos:} coluna inteira ($\sim$28~m) ou análise plug a
            plug nos dez pontos com imagem 3D.
    \end{itemize}
    \vspace{0.02cm}
    {\scriptsize
    \emph{Perfil enriquecido}: 87 linhas com perfil e lab em todo o intervalo;
    CT só onde há plug. Nas 87 linhas, $X$ = curvas de perfil; $y$ = lab.}
    \column{0.48\textwidth}
    \centering
    \begin{tikzpicture}[scale=0.68, every node/.style={font=\scriptsize}]
      \draw[fill=po861blue!12, draw=po861blue, thick, rounded corners=2pt]
        (-0.2,5.6) rectangle (4.6,6.5);
      \node[align=center] at (2.2,6.05)
        {\textbf{Perfil + laboratório}\\87 profundidades};
      \draw[fill=po861green!12, draw=po861green, thick, rounded corners=2pt]
        (-0.2,3.7) rectangle (4.6,4.6);
      \node[align=center] at (2.2,4.15)
        {\textbf{MicroCT}\\10 plugs};
      \draw[fill=orange!12, draw=orange!80!black, thick, rounded corners=2pt]
        (-0.2,1.8) rectangle (4.6,2.7);
      \node[align=center] at (2.2,2.25)
        {\textbf{Alinhamento $\pm$0,5~m}\\8 prof.\ com CT};
      \draw[fill=po861blue!8, draw=po861blue, thick, rounded corners=2pt]
        (-0.2,0.0) rectangle (4.6,1.0);
      \node[align=center] at (2.2,0.5)
        {\textbf{Perfil enriquecido} ($n=87$)};
      \draw[->, thick, po861blue] (2.2,3.6) -- (2.2,2.8);
      \draw[->, thick, po861blue] (2.2,5.5) -- (2.2,4.7);
      \draw[->, thick, po861blue] (2.2,1.7) -- (2.2,1.1);
    \end{tikzpicture}
  \end{columns}
\end{frame}

\begin{frame}{Entradas ($X$) e alvos ($y$) dos experimentos}
  \footnotesize
  Cada modelo prevê \textbf{um} alvo de laboratório a partir de
  \textbf{oito} curvas de perfil. Lab fora de $X$ para evitar vazamento.
  \vspace{0.06cm}
  \setlength{\columnsep}{0.05\textwidth}
  \begin{columns}[T,onlytextwidth]
    \column{0.45\textwidth}
    \begin{block}{Features ($X$)}
      \scriptsize
      Oito curvas medidas \emph{in situ} (87 profundidades):
      \vspace{0.04cm}
      \centering
      \setlength{\tabcolsep}{3pt}
      \renewcommand{\arraystretch}{1.05}
      \begin{tabular}{@{}>{\raggedright\arraybackslash}p{1.65cm}>{\raggedright\arraybackslash}p{2.45cm}@{}}
        \toprule
        Variável & Significado \\
        \midrule
        GR, $\rho_b$ & litologia / densidade \\
        $R_t$ profunda e rasa & resistividade \\
        $\phi_{\mathrm{N}}$, $\phi_{\mathrm{S}}$, \phiND{} & porosidade estimada \\
        Litotipo & cod.\ 1--4 (legenda abaixo) \\
        \bottomrule
      \end{tabular}
      \vspace{0.04cm}

      {\scriptsize Total: 8 features wireline em cada linha.}
    \end{block}
    \column{0.45\textwidth}
    \begin{block}{Alvos ($y$)}
      \scriptsize
      Grandezas de laboratório na mesma profundidade:
      \vspace{0.04cm}
      \centering
      \setlength{\tabcolsep}{3pt}
      \renewcommand{\arraystretch}{1.05}
      \begin{tabular}{@{}>{\raggedright\arraybackslash}p{1.0cm}>{\raggedright\arraybackslash}p{0.9cm}>{\raggedright\arraybackslash}p{2.2cm}@{}}
        \toprule
        Alvo & Tipo & Papel \\
        \midrule
        \philab & regressão & porosidade medida \\
        \fzilab & regressão & qualid.\ zona de fluxo \\
        \klab & regressão & permeabilidade \\
        HFU (1--4) & classif. & classe de fluxo \\
        \bottomrule
      \end{tabular}
      \vspace{0.04cm}

      {\scriptsize HFU: 42 / 29 / 10 / 6 amostras/classe.}
    \end{block}
  \end{columns}
  \vspace{0.06cm}
  {\scriptsize
  \textbf{Dois datasets:} perfil enriquecido ($n=87$, 5205,9--5233,7~m);
  dez plugs com microCT (CT só no segundo caso).\\
  Litotipos:\\
  1=Grainstone, 2=Shrubs, 3=Laminite, 4=Spherulite.}
\end{frame}

% ------------------------------------------------------------------
\section{Protocolo}

\begin{frame}{Dois protocolos de validação}
  \vspace{-0.25cm}
  {\footnotesize
  $y$ = medidas de laboratório; lab nunca em $X$.
  Modelos testados: RF, GB, XGB, MLP, LR.}
  \vspace{0.06cm}

  \begin{columns}[T,onlytextwidth]
    \column{0.49\textwidth}
    \begin{block}{Cenário A --- perfil completo ($n=87$)}
      \scriptsize
      \textbf{$X$:} oito curvas wireline \quad
      \textbf{$y$:} \philab, \fzilab, \klab, HFU \quad
      \textbf{CV:} 3 blocos OOF
      \vspace{0.05cm}

      \centering
      \begin{tikzpicture}[
        node distance=0.35cm and 0.22cm,
        mlbox/.style={
          draw=po861blue, thick, rounded corners=2pt, fill=po861blue!10,
          minimum height=0.62cm, minimum width=1.45cm,
          align=center, font=\scriptsize, inner sep=1pt},
        mlarr/.style={-{Stealth[length=2mm]}, thick, po861blue},
      ]
        \node[mlbox] (xa) {$X$\\8 curvas};
        \node[mlbox, right=of xa] (cva) {3 blocos\\OOF};
        \node[mlbox, right=of cva] (ma) {$R^2$, RMSE};
        \draw[mlarr] (xa) -- (cva) -- (ma);
        \node[font=\tiny, text=po861blue, below=0.12cm of cva]
          {treino 2 $\cdot$ teste 1 bloco};
      \end{tikzpicture}
      \vspace{0.06cm}

      \centering
      \begin{tabular}{ccrr}
        \toprule
        Bloco & Prof.\ (m) & $\overline{\phi}$ (pu) & $\overline{k}$ (mD) \\
        \midrule
        0 & 5206--5214 & 0,092 & 3,6 \\
        1 & 5215--5223 & 0,115 & 9,5 \\
        2 & 5225--5234 & 0,145 & 54,6 \\
        \bottomrule
      \end{tabular}
    \end{block}

    \column{0.49\textwidth}
    \begin{exampleblock}{Cenário B --- plugs + microCT ($n=10$)}
      \scriptsize
      \textbf{$X$:} oito curvas + descritores CT \quad
      \textbf{$y$:} \philab{} ou \fzilab \quad
      \textbf{CV:} LOPO
      \vspace{0.05cm}

      \centering
      \begin{tikzpicture}[
        node distance=0.35cm and 0.22cm,
        mlbox/.style={
          draw=po861green!70!black, thick, rounded corners=2pt,
          fill=po861green!12, minimum height=0.62cm, minimum width=1.45cm,
          align=center, font=\scriptsize, inner sep=1pt},
        mlarr/.style={
          -{Stealth[length=2mm]}, thick, po861green!60!black},
      ]
        \node[mlbox] (xb) {$X$\\8 curvas\\+ CT};
        \node[mlbox, right=of xb] (cvb) {LOPO\\1 plug};
        \node[mlbox, right=of cvb] (mb) {$R^2$, RMSE};
        \draw[mlarr] (xb) -- (cvb) -- (mb);
        \node[font=\tiny, text=po861green!60!black, below=0.12cm of cvb]
          {treino 9 $\cdot$ teste 1 plug};
      \end{tikzpicture}
      \vspace{0.06cm}

      {\tiny
      CT entra só nos plugs, não nas 87 linhas.
      Comparamos RF com e sem descritores CT (LOPO): CT não ajudou
      em \philab{} nem \fzilab{} (próximo slide).}
    \end{exampleblock}
  \end{columns}

  \vspace{0.04cm}
  {\scriptsize
  Sem vazamento: laboratório não prediz laboratório.
  O número que vale é o OOF: o bloco (ou o plug) avaliado não entrou
  no treino. Sorteio 80/20 infla o $R^2$ neste tipo de dado.}

  \vspace{0.04cm}
  \begin{block}{Falha no bloco~2 (Cenário~A)}
    \scriptsize
    Entre 5225--5234~m a coluna fica mais permeável ($\overline{k}=54{,}6$~mD)
    e muda de litologia: some o Grainstone (litotipo~1); predominam Shrubs
    e Spherulite. O RF de \philab{} treinado em cima não segura essa faixa
    ($R^2<0$ no teste OOF).
  \end{block}
\end{frame}

% ------------------------------------------------------------------
\section{Diagnóstico}

\begin{frame}{Porosidade sim, FZI quase não}
  \begin{columns}[T,onlytextwidth]
    \column{0.44\textwidth}
    \footnotesize
    \textbf{$|r|$ com \fzilab}
    \vspace{0.04cm}
    \begin{tabular}{lr}
      \toprule
      Variável & $|r|$ \\
      \midrule
      RQI (lab) & 0,887 \\
      \klab{} (lab) & 0,709 \\
      HFU (lab) & 0,614 \\
      \phiND{} (perfil) & 0,308 \\
      GR (perfil) & 0,040 \\
      \bottomrule
    \end{tabular}
    \vspace{0.10cm}
    \textbf{Com \philab:} $\phi_{\mathrm{N}}$ ($r \approx 0{,}74$),
    \phiND{} ($r \approx 0{,}69$).

    \vspace{0.08cm}
    {\scriptsize FZI depende de textura e $k$ --- o perfil mal captura isso:
    o melhor $r$ de curva (\phiND) fica em 0,31.}
    \column{0.54\textwidth}
    \centering
    \slidefig[0.72]{log_features_vs_FZI_lab_corr}
    \slidecaption{Correlação de Pearson: curvas de perfil versus \fzilab}
  \end{columns}
\end{frame}

\begin{frame}{Teste \emph{oracle}: FZI precisa de lab, não de perfil}
  \begin{columns}[T,onlytextwidth]
    \column{0.50\textwidth}
    \footnotesize
    \begin{itemize}
      \item RF só com perfis: \fzilab{} falha (bloco~1, $R^2 \approx -1{,}95$).
      \item \emph{Oracle} (diagnóstico): entram \philab{} e \klab{} ---
            informação que só teríamos depois do testemunho.
      \item Com isso, $R^2$ sobe para 0,37--0,88 nos blocos~0 e~1
            (média $\approx 0{,}69$).
    \end{itemize}
    \vspace{0.06cm}
    {\scriptsize FZI tem sinal previsível, mas vem de porosidade e $k$ de
    laboratório --- não das oito curvas de perfil. O problema não é o
    algoritmo.}
    \column{0.48\textwidth}
    \centering
    \slidefig[0.80]{FZI_oracle_vs_logs_r2_by_fold}
    \slidecaption{$R^2$ por bloco: perfil isolado versus \emph{oracle}}
  \end{columns}
\end{frame}

% ------------------------------------------------------------------
\section{Resultados no perfil}

\begin{frame}{RF nas 87 linhas: só \philab{} com $R^2$ OOF positivo}
  \footnotesize
  Mesmo RF, mesmo protocolo (\oofcv): oito curvas $\rightarrow$ um alvo.
  Cada linha da tabela é um experimento separado.
  $R^2<0$ = pior que prever a média do bloco de teste.
  \vspace{0.06cm}
  \begin{columns}[T,onlytextwidth]
    \column{0.52\textwidth}
    \begin{tabular}{lrrr}
      \toprule
      Alvo ($y$) & RMSE & \rSqOOF & std \\
      \midrule
      \philab{} (pu) & 0,033 & \textbf{0,146} & 0,222 \\
      \fzilab & 1,813 & $-0{,}819$ & 0,820 \\
      \klab{} (mD) & 76,2 & $-15{,}4$ & 18,4 \\
      \bottomrule
    \end{tabular}
    \vspace{0.06cm}
    {\scriptsize
    Só \philab{} serve como porosidade preliminar de perfil.
    FZI e $k$ negativos não indicam RF ruim em \philab{} --- são alvos sem
    sinal wireline.}
    \column{0.46\textwidth}
    \centering
    \slidefig[0.82]{target_comparison_mean_r2}
    \slidecaption{\rSqOOF{} por alvo (RF); referência = média do treino}
  \end{columns}
  \vspace{0.04cm}
  {\footnotesize $R^2 \approx 0{,}15$ em \philab{} é modesto, mas positivo ---
  dá para usar como entrada de DEM/SC, sem substituir testemunho.}
\end{frame}

\begin{frame}{Cinco regressores em \philab{}: só RF acima de $R^2=0$}
  \footnotesize
  Alvo fixo: \philab{} ($n=87$, \oofcv). Comparação de algoritmos
  \textbf{só para porosidade} --- FZI e $k$ ficam no slide anterior.
  \vspace{0.06cm}
  \begin{columns}[T,onlytextwidth]
    \column{0.58\textwidth}
    \textbf{\philab{} (pu) --- comparação entre algoritmos}
    \vspace{0.04cm}
    \begin{tabular}{lrr}
      \toprule
      Algoritmo & RMSE & \rSqOOF \\
      \midrule
      Random Forest & 0,033 & \textbf{0,146} \\
      Regressão linear & 0,036 & $-0{,}261$ \\
      XGBoost & 0,036 & $-0{,}099$ \\
      Grad.\ Boosting & 0,040 & $-0{,}313$ \\
      MLP & 0,097 & $-6{,}871$ \\
      \bottomrule
    \end{tabular}
    \vspace{0.08cm}
    {\scriptsize
    Nos outros alvos (mesma validação): FZI $R^2_{\mathrm{RF}}=-0{,}82$;
    $k$ $R^2_{\mathrm{RF}}=-15{,}4$. Nenhum regressor se salva.}
    \column{0.40\textwidth}
    \begin{itemize}
      \item Em \philab, só RF fica com $R^2>0$.
      \item Trocar algoritmo não compensa falta de informação no perfil.
      \item Próximo slide: \phiND{} direto versus RF em \philab.
    \end{itemize}
  \end{columns}
\end{frame}

\begin{frame}{\phiND{} direto supera o RF para \philab{}}
  \vspace{-0.28cm}
  \begin{columns}[T,onlytextwidth]
    \column{0.31\textwidth}
    \footnotesize
    A porosidade neutrão--densidade (\phiND{}) funciona como proxy
    direta de \philab{} --- sem aprendizado de máquina.
    \vspace{0.04cm}

    \centering
    \setlength{\tabcolsep}{4pt}
    \scriptsize
    \begin{tabular}{lrrr}
      \toprule
      Método & RMSE & MAE & $R^2$ \\
      \midrule
      \phiND{} direto & 0,032 & 0,022 & \textbf{0,446} \\
      RF OOF & 0,036 & 0,026 & 0,300 \\
      \bottomrule
    \end{tabular}
    \vspace{0.05cm}

    {\scriptsize
    $r(\phiND{}, \philab) = 0{,}701$ ($n=87$).\\[0.06cm]
    \textbf{Acima:} \philab{} (pontos), \phiND{} (azul), RF (verde).\\[0.03cm]
    \textbf{Abaixo:} métricas e dispersão 1:1.\\[0.03cm]
    \textbf{Conclusão:} não compensa substituir \phiND{} por RF.}
    \column{0.67\textwidth}
    \centering
    \slidefig[0.58]{phi_nd_vs_rf_depth_profile}
  \end{columns}
  \vspace{-0.12cm}
  \centering
  \slidefig[0.36]{phi_nd_vs_rf_bottom_panels}
\end{frame}

\begin{frame}{\philab{} RF: métricas OOF por bloco de profundidade}
  \begin{columns}[T,onlytextwidth]
    \column{0.44\textwidth}
    \footnotesize
    \textbf{Random Forest, alvo \philab{}}
    \vspace{0.04cm}
    \begin{tabular}{ccrr}
      \toprule
      Bloco & Prof.\ (m) & RMSE & $R^2$ OOF \\
      \midrule
      0 & 5206--5214 & 0,015 & 0,367 \\
      1 & 5215--5223 & 0,037 & 0,227 \\
      2 & 5225--5234 & 0,047 & $-0{,}157$ \\
      \midrule
      Média & -- & 0,033 & 0,146 \\
      \bottomrule
    \end{tabular}
    \vspace{0.08cm}
    {\scriptsize
    $R^2$ negativo só no bloco~2 (faixa mais profunda e heterogênea).
    A média dos três blocos segue positiva (0,146).
    Sorteio 80/20 inflaria o $R^2$ --- por isso não o usamos.}
    \column{0.54\textwidth}
    \centering
    \slidefig[0.78]{phi_lab_oof_rmse_by_depth_fold}
    \slidecaption{RMSE OOF por bloco: RF, GB, MLP e LR (\philab{})}
  \end{columns}
\end{frame}

\begin{frame}{\philab{} em profundidade: predições OOF}
  \footnotesize
  Mesmo protocolo das tabelas: $n=87$, 8 wireline, \oofcv.
  Linhas tracejadas vermelhas = limites entre blocos de teste.
  \vspace{0.04cm}
  \begin{columns}[T,onlytextwidth]
    \column{0.56\textwidth}
    \centering
    \slidefig[0.88]{phi_lab_oof_depth_panels}
    \slidecaption{Observado (tracejado) versus OOF (colorido) por regressor}
    \column{0.42\textwidth}
    \begin{itemize}
      \item \textbf{RF (verde):} acompanha \philab{} nos blocos 0--1;
            perde no bloco profundo.
      \item \textbf{LR (azul):} próximo do RF neste poço.
      \item \textbf{MLP (vermelho):} instável ($R^2 \ll 0$).
      \item O erro cresce com a profundidade em todos os modelos.
    \end{itemize}
  \end{columns}
\end{frame}

% ------------------------------------------------------------------
\section{Limites e alternativas}

\begin{frame}{MicroCT nos 10 plugs: RF não ganhou com CT}
  \begin{columns}[T,onlytextwidth]
    \column{0.44\textwidth}
    \footnotesize
    \textbf{$R^2$ global LOPO}
    \vspace{0.04cm}
    \begin{tabular}{lrr}
      \toprule
      Alvo & Só perfil & Perfil + CT \\
      \midrule
      \fzilab & 0,151 & 0,116 \\
      \philab & $-0{,}334$ & $-0{,}466$ \\
      \bottomrule
    \end{tabular}
    \vspace{0.10cm}
    \begin{itemize}
      \item Com dez pontos, qualquer $R^2$ oscila bastante.
      \item Incluir CT \emph{piorou} o RF --- provável overfitting.
      \item O CT volta na Etapa~2, para geometria no DEM/SC.
    \end{itemize}
    \column{0.54\textwidth}
    \centering
    \slidefig[0.76]{plug_out_pred_vs_obs_rf}
    \slidecaption{FZI nos plugs: predito versus observado (LOPO, só perfil)}
  \end{columns}
\end{frame}

\begin{frame}{HFU pelo perfil: resultado fraco}
  \begin{columns}[T,onlytextwidth]
    \column{0.46\textwidth}
    \footnotesize
    \textbf{Classificação OOF (blocos)}
    \vspace{0.04cm}
    \begin{tabular}{lrrr}
      \toprule
      Modelo & Acc. & Acc.\ bal. & F1 \\
      \midrule
      Logística bal. & 0,356 & 0,220 & 0,217 \\
      RF balanceado & 0,299 & 0,179 & 0,161 \\
      \bottomrule
    \end{tabular}
    \vspace{0.08cm}
    {\scriptsize Quatro classes: acaso $\approx 25\%$. Sempre chutar HFU1
    daria 48\% sem aprender nada.}

    \vspace{0.08cm}
    {\scriptsize Na Etapa~2 usamos HFU de laboratório; o classificador fica
    só como recurso onde lab faltar.}
    \column{0.52\textwidth}
    \centering
    \slidefig[0.78]{hfu_model_comparison_oof}
    \slidecaption{Métricas OOF: logística balanceada versus RF balanceado}
  \end{columns}
\end{frame}

\begin{frame}{GAM e Ridge não batem o RF em \philab}
  \begin{columns}[T,onlytextwidth]
    \column{0.44\textwidth}
    \footnotesize
    \begin{tabular}{lrr}
      \toprule
      Modelo & RMSE (pu) & $R^2$ \\
      \midrule
      RF & 0,033 & \textbf{0,146} \\
      GAM (spline+Ridge) & 0,036 & $-0{,}052$ \\
      Ridge & 0,037 & $-0{,}261$ \\
      \bottomrule
    \end{tabular}
    \vspace{0.10cm}
    Com 87 amostras e oito curvas, modelos lineares e splines não seguram
    a mudança litológica entre blocos.
    \column{0.54\textwidth}
    \centering
    \slidefig[0.78]{phi_model_comparison_mean_r2}
    \slidecaption{$R^2$ médio: RF versus GAM e Ridge}
  \end{columns}
\end{frame}

% ------------------------------------------------------------------
\section{Encaminhamento}

\begin{frame}{Da Etapa~1 à física de rochas}
  \vspace{-0.12cm}
  \begin{columns}[T,onlytextwidth]
    \column{0.50\textwidth}
    \footnotesize
    \textbf{Fechamos na Etapa~1}
    \begin{itemize}
      \item Porosidade de entrada: \phiND{} do perfil (RF como reserva).
      \item FZI e $k$: fora do escopo com perfil isolado.
      \item HFU de laboratório para fatiar a modelagem física.
    \end{itemize}
    \textbf{Próximo}
    \begin{itemize}
      \item Etapa~2: DEM/SC e Gassmann por HFU, microCT nos dez plugs.
      \item Etapa~3: ML no resíduo de $V_p$, depois de validar a física.
    \end{itemize}
    \column{0.48\textwidth}
    \centering
    \begin{tikzpicture}[
      node distance=0.95cm,
      flowbox/.style={
        draw, thick, rounded corners=4pt,
        align=center, text width=5.5cm,
        inner sep=5pt, font=\scriptsize
      },
      flowarr/.style={-{Stealth[length=2.8mm]}, thick, po861blue},
      arrlabel/.style={
        midway, fill=white, inner sep=2pt,
        font=\tiny, text=po861blue, align=center
      }
    ]
      \node[flowbox, fill=po861blue!18, draw=po861blue] (e1) {%
        \textbf{1. Diagnóstico ML}\\
        Limites do perfil
      };
      \node[flowbox, fill=po861green!18, draw=po861green, below=of e1] (e2) {%
        \textbf{2. DEM/SC + Gassmann}\\
        microCT + HFU $\rightarrow$ $V_p$, $V_s$
      };
      \node[flowbox, fill=orange!18, draw=orange!70!black, below=of e2] (e3) {%
        \textbf{3. Física + ML residual}\\
        Ajuste fino de $V_p$ (RF / Ridge OOF)
      };
      \draw[flowarr] (e1.south) --
        node[arrlabel] {porosidade, HFU} (e2.north);
      \draw[flowarr] (e2.south) --
        node[arrlabel] {resíduo} (e3.north);
    \end{tikzpicture}
  \end{columns}
\end{frame}

\begin{frame}{Conclusões}
  \footnotesize
  \begin{columns}[T,onlytextwidth]
    \column{0.56\textwidth}
    \begin{enumerate}
      \item \phiND{} direto como porosidade para a física ($R^2 = 0{,}45$;
            RF OOF fica em 0,30).
      \item FZI só com perfil: não usar como modelo.
      \item HFU de laboratório na Etapa~2; classificador automático em segundo plano.
      \item microCT para DEM/SC, não para enriquecer RF de perfil.
      \item Física de rochas antes de mais \emph{tuning} perfil $\rightarrow$ $V_p$.
    \end{enumerate}
    \column{0.40\textwidth}
    \centering
    \textbf{Referência rápida}\\[0.08cm]
    \begin{tabular}{lc}
      \toprule
      Grandeza & $R^2$ ou acc. \\
      \midrule
      \philab & 0,146 \\
      \fzilab & $-0{,}819$ \\
      \klab & $-15{,}4$ \\
      HFU & 36\% \\
      \bottomrule
    \end{tabular}
  \end{columns}
  \vspace{0.12cm}
  \noindent\footnotesize O perfil ajuda na porosidade via \phiND{}; RF nao
  supera essa proxy. Para textura e permeabilidade, seguimos com DEM/SC e HFU
  de laboratório na Etapa~2.
  \vspace{0.15cm}
  \begin{center}
    {\small Obrigado.}
  \end{center}
\end{frame}

\end{document}
